\section{Техническое задание}

Проект представляет собой веб-приложение <<Цветовая доска>>, предназначенное для анализа цветов, генерации оттенков исходного цвета и определения контрастности текста относительно фонового цвета.

Приложение позволяет пользователю добавлять один или несколько цветов, получать для каждого цвета набор оттенков по яркости, рассчитывать контрастность текста на выбранном фоне, а также переводить цвет между различными цветовыми форматами.

\subsection*{Функциональные возможности}

Приложение должно предоставлять следующие возможности:

\begin{itemize}
    \item добавление одного или нескольких исходных цветов;
    \item ввод цвета в различных форматах: HEX, sRGB, LCH, OKLCH, CMYK;
    \item автоматическое определение и проверка корректности введённого цвета;
    \item перевод цвета из одного формата в другой;
    \item выбор формата отображения цвета через переключатель или выпадающий список;
    \item разбиение исходного цвета на набор оттенков по яркости;
    \item отображение палитры оттенков для каждого добавленного цвета;
    \item расчёт контраста между цветом текста и цветом фона;
    \item автоматический подбор подходящего цвета текста: белого или чёрного;
    \item отображение результата проверки контрастности по требованиям WCAG;
    \item удаление добавленного цвета;
    \item копирование значения цвета в выбранном формате.
\end{itemize}

Расчёт контраста выполняется по формуле WCAG, представленной в спецификации Web Content Accessibility Guidelines в определении contrast ratio:
\texttt{https://www.w3.org/TR/WCAG21/\#dfn-contrast-ratio}

\[
K = \frac{L_1 + 0.05}{L_2 + 0.05},
\]

где \(L_1\)~--- относительная яркость более светлого цвета, \(L_2\)~--- относительная яркость более тёмного цвета. Значение контраста находится в диапазоне от \(1:1\) до \(21:1\).

Для обычного текста минимальное рекомендуемое значение контраста составляет \(4.5:1\), для крупного текста~--- \(3:1\).

\subsection*{Вводимые пользователем данные}

Пользователь может вводить следующие данные:

\begin{itemize}
    \item исходный цвет;
    \item формат исходного цвета;
    \item формат, в который необходимо перевести цвет;
    \item количество оттенков в палитре;
    \item цвет текста для ручной проверки контраста;
    \item цвет фона;
    \item название цвета или палитры.
\end{itemize}

Примеры допустимых форматов цвета:

\begin{verbatim}
HEX:   #3B82F6
sRGB:  rgb(59, 130, 246)
LCH:   lch(60% 70 250)
OKLCH: oklch(65% 0.18 250)
CMYK:  cmyk(76% 47% 0% 4%)
\end{verbatim}

\subsection*{Требования к обработке цвета}

Система должна выполнять следующие операции:

\begin{itemize}
    \item распознавать формат введённого цвета;
    \item проверять корректность значений компонентов цвета;
    \item приводить цвет к внутреннему универсальному представлению;
    \item выполнять конвертацию между поддерживаемыми форматами;
    \item отображать результат конвертации без потери читаемости;
    \item корректно обрабатывать цвета с высокой и низкой яркостью;
    \item корректно работать с несколькими цветами одновременно.
\end{itemize}

Для расчёта контраста цвета приводятся к цветовому пространству sRGB, так как формула WCAG использует относительную яркость, вычисляемую на основе sRGB-компонентов.

Относительная яркость вычисляется по формуле:

\[
L = 0.2126R + 0.7152G + 0.0722B,
\]

где \(R\), \(G\), \(B\)~--- линейные значения красного, зелёного и синего каналов.

Для перевода компонента sRGB в линейное значение используется формула:

\[
C_{\text{linear}} =
\begin{cases}
\frac{C_{\text{srgb}}}{12.92}, & C_{\text{srgb}} \leq 0.04045,\\
\left(\frac{C_{\text{srgb}} + 0.055}{1.055}\right)^{2.4}, & C_{\text{srgb}} > 0.04045.
\end{cases}
\]

\subsection*{Разбиение цвета на оттенки}

Для каждого добавленного цвета приложение должно строить палитру оттенков по яркости.

Разбиение цвета выполняется в цветовом пространстве OKLCH или LCH, так как эти пространства позволяют удобно изменять воспринимаемую яркость цвета. При генерации оттенков изменяется компонент яркости, а тон и насыщенность по возможности сохраняются.

Пример набора оттенков:

\begin{verbatim}
50  -- очень светлый оттенок
100 -- светлый оттенок
200
300
400
500 -- исходный или близкий к исходному цвет
600
700
800
900 -- тёмный оттенок
950 -- очень тёмный оттенок
\end{verbatim}

Для каждого оттенка необходимо отображать:

\begin{itemize}
    \item визуальный образец цвета;
    \item значение цвета в выбранном формате;
    \item рекомендуемый цвет текста;
    \item коэффициент контраста с белым текстом;
    \item коэффициент контраста с чёрным текстом;
    \item статус соответствия требованиям WCAG.
\end{itemize}

\subsection*{Требования к контрасту}

Для каждого цвета или оттенка система должна рассчитывать контраст:

\begin{itemize}
    \item с белым текстом;
    \item с чёрным текстом;
    \item с пользовательским цветом текста, если он задан вручную.
\end{itemize}

Результат проверки контраста должен отображаться в понятном виде:

\begin{verbatim}
Контраст с белым: 5.82:1 -- AA
Контраст с чёрным: 3.61:1 -- не подходит для обычного текста
Рекомендуемый текст: белый
\end{verbatim}

Критерии оценки контрастности:

\begin{itemize}
    \item \(K \geq 7\)~--- соответствует уровню AAA для обычного текста;
    \item \(K \geq 4.5\)~--- соответствует уровню AA для обычного текста;
    \item \(K \geq 3\)~--- подходит для крупного текста;
    \item \(K < 3\)~--- недостаточный контраст.
\end{itemize}

\subsection*{Интерфейс приложения}

Интерфейс приложения должен быть выполнен в виде доски.

Основные области интерфейса:

\begin{enumerate}
    \item \textbf{Панель ввода цвета.}

    Содержит поле ввода цвета, выбор формата, кнопку добавления цвета и настройки генерации оттенков.

    \item \textbf{Доска цветов.}

    Содержит карточки всех добавленных цветов. Каждая карточка показывает исходный цвет, его значения в разных форматах и палитру оттенков.

    \item \textbf{Блок контраста.}

    Для каждого оттенка отображается результат проверки контрастности с белым, чёрным и пользовательским цветом текста.

    \item \textbf{Блок конвертации.}

    Пользователь может выбрать формат отображения цвета через выпадающий список или переключатель. После выбора все цвета и оттенки отображаются в выбранном формате.
\end{enumerate}

Пример структуры карточки цвета:

\begin{verbatim}
Цвет: #3B82F6

sRGB:  rgb(59, 130, 246)
OKLCH: oklch(...)
LCH:   lch(...)
CMYK:  cmyk(...)

Оттенок 500:
Фон: #3B82F6
Текст: #FFFFFF
Контраст: 5.82:1
Статус: AA
\end{verbatim}

\subsection*{Технологический стек}

Frontend-часть приложения реализуется с использованием Vue 3.

Рекомендуемый стек frontend-части:

\begin{itemize}
    \item Vue 3;
    \item TypeScript;
    \item Vite;
    \item Pinia для хранения состояния;
    \item CSS, SCSS или Tailwind CSS для оформления интерфейса.
\end{itemize}

Backend-часть приложения реализуется на языке Python с использованием FastAPI.

Для работы с цветами используется библиотека ColorAide. Она применяется для парсинга введённых пользователем цветов, конвертации между цветовыми пространствами, работы с форматами sRGB, LCH, OKLCH, Lab и генерации оттенков через изменение компонента яркости.

Расчёт контрастности выполняется отдельно по формуле WCAG на основе относительной яркости sRGB.

\subsection*{Требования к данным}

Ограничения на вводимые данные:

\begin{itemize}
    \item длина строки цвета~--- до 100 символов;
    \item количество одновременно добавленных цветов~--- до 20;
    \item количество оттенков для одного цвета~--- от 3 до 20;
    \item поддерживаемые форматы: HEX, sRGB, LCH, OKLCH, CMYK;
    \item некорректные значения должны сопровождаться понятным сообщением об ошибке.
\end{itemize}

Примеры сообщений об ошибках:

\begin{verbatim}
Некорректный HEX-цвет.
Компонент RGB должен быть в диапазоне от 0 до 255.
Значение яркости в OKLCH должно быть от 0% до 100%.
Невозможно перевести цвет в выбранный формат.
\end{verbatim}

\subsection*{Требования к хранению состояния}

Приложение должно хранить:

\begin{itemize}
    \item список добавленных цветов;
    \item выбранный формат отображения;
    \item настройки генерации оттенков;
    \item пользовательский цвет текста для проверки контраста.
\end{itemize}

Состояние приложения хранится во frontend-части с использованием Pinia. Дополнительно допускается сохранение данных в \texttt{localStorage}, чтобы после перезагрузки страницы доска не очищалась.

\subsection*{Требования к дизайну}

Интерфейс должен быть минималистичным и удобным для сравнения цветов.

Требования к интерфейсу:

\begin{itemize}
    \item карточки цветов должны хорошо читаться;
    \item оттенки должны отображаться в виде горизонтальной палитры или сетки;
    \item для каждого оттенка должен быть виден рекомендуемый цвет текста;
    \item успешный и неуспешный результат проверки контраста должны визуально отличаться;
    \item приложение должно корректно работать на экранах ноутбуков и настольных компьютеров;
    \item интерфейс должен поддерживать светлую и тёмную тему.
\end{itemize}

\subsection*{Минимальный сценарий работы пользователя}

\begin{enumerate}
    \item Пользователь открывает приложение.
    \item Вводит цвет, например \texttt{\#3B82F6}.
    \item Выбирает формат исходного цвета.
    \item Нажимает кнопку <<Добавить цвет>>.
    \item Приложение создаёт карточку цвета.
    \item Для цвета автоматически строится палитра оттенков.
    \item Для каждого оттенка рассчитывается контраст с белым и чёрным текстом.
    \item Пользователь выбирает формат отображения, например OKLCH.
    \item Все цвета на доске отображаются в выбранном формате.
    \item Пользователь добавляет ещё несколько цветов и сравнивает их палитры.
\end{enumerate}

\subsection*{Планируемый результат}

В результате должен быть реализован веб-сервис, позволяющий дизайнеру или разработчику быстро анализировать цвета, получать оттенки исходного цвета, проверять читаемость текста на фоне и переводить цвета между популярными цветовыми форматами.

Приложение может использоваться при разработке пользовательских интерфейсов, дизайн-систем, цветовых палитр и проверке доступности веб-страниц.
